\section{Introduzione e obiettivi}
\label{sec:introduzione}
\thispagestyle{plain}
{\itshape In questo primo capitolo viene presentato il contesto applicativo e metodologico alla base del progetto, definendo le sfide di sicurezza tipiche della gestione di database distribuiti e illustrando i principi teorici del paradigma Zero Trust. Si introduce quindi il caso di studio ospedaliero scelto come dominio di riferimento e si delineano gli obiettivi e la struttura complessiva della relazione.}

\subsection{La protezione dei dati aziendali e delle risorse critiche}
Le infrastrutture IT aziendali rappresentano uno dei bersagli più ambiti e critici per gli attori malevoli a livello globale. La digitalizzazione dei servizi e la transizione verso sistemi in cloud hanno introdotto enormi benefici in termini di efficienza operativa, tra cui l'accesso condiviso a risorse sensibili, l'automazione dei flussi e il monitoraggio IoT di asset industriali. Tuttavia, questa evoluzione ha contemporaneamente esteso a dismisura la superficie d'attacco delle organizzazioni.

La criticità di questo scenario è dettata da molteplici fattori. In primo luogo, vi è la natura estremamente riservata delle informazioni trattate all'interno dei database, che includono dati finanziari, proprietà intellettuali, credenziali e profili di utenti. La perdita di riservatezza o l'alterazione di tali dati viola gravemente le principali normative vigenti sulla protezione dei dati personali e sulla sicurezza delle reti a livello internazionale. In secondo luogo, va considerato l'impatto sulla \textbf{business continuity}: a differenza dei semplici disservizi temporanei, un attacco informatico di tipo ransomware che cripta o altera i database operativi può causare il blocco completo dei servizi, enormi perdite finanziarie e danni reputazionali incalcolabili. Infine, l'eterogeneità degli accessi complica ulteriormente il quadro: gli operatori e gli amministratori devono poter consultare i dati sia dai terminali interni alla rete aziendale, come le workstation d'ufficio, sia da remoto, ad esempio in regime di smart working, introducendo vettori di minaccia esterni non controllati.

\subsection{La filosofia Zero Trust: ``never Trust, always Verify''}
La sicurezza perimetrale classica, basata sul modello a ``castello'', in cui un firewall protegge una rete interna considerata intrinsecamente sicura da una rete esterna ostile, si è dimostrata obsoleta e inefficace. Una volta che un attaccante riesce a penetrare il perimetro, tramite phishing, exploit o \textbf{insider threat}, ottiene totale libertà di movimento laterale.

Il paradigma \textbf{Zero Trust Architecture} (ZTA), codificato dal NIST nel documento speciale \textit{SP 800-207}, supera questa concezione abbattendo il concetto di fiducia implicita basata sulla localizzazione di rete. I cardini della filosofia ZTA sono tre:
\begin{enumerate}
    \item \textbf{Mai fidarsi, verificare sempre (Never Trust, Always Verify):} ogni singola richiesta di accesso ad una risorsa, indipendentemente dalla rete da cui proviene, sia essa interna o esterna, deve essere autenticata, autorizzata e validata crittograficamente prima dell'inoltro.
    \item \textbf{Accesso con privilegi minimi (Least Privilege):} l'accesso viene concesso in modo estremamente granulare e ristretto solo alle risorse necessarie per lo svolgimento della specifica attività, ad esempio limitando la scrittura e la cancellazione di record sensibili ai soli operatori autorizzati.
    \item \textbf{Presumere la violazione (Assume Breach):} progettare il sistema assumendo che alcuni segmenti siano già compromessi. Ciò impone l'adozione di microsegmentazione di rete per limitare il raggio d'azione dell'attaccante, crittografia end-to-end per proteggere i dati in transito e sistemi di monitoraggio continuo del traffico.
\end{enumerate}

\subsection{Il caso di studio: l'infrastruttura ospedaliera}
Il dominio applicativo scelto per la validazione dell'architettura è quello di una \textbf{struttura ospedaliera}, un contesto in cui la criticità dei dati trattati e l'eterogeneità degli accessi rendono particolarmente evidente la necessità di un approccio Zero Trust.

Il sistema gestisce un database MongoDB contenente un esempio di informazioni sanitarie altamente sensibili, come anagrafiche dei pazienti, cartelle cliniche e dati diagnostici. L'accesso a queste risorse è richiesto da tre categorie di operatori con profili di rischio differenti:
\begin{itemize}
    \item \textbf{Alice}: personale medico interno all'ospedale, dotato di workstation aziendale con chip \textbf{TPM} hardware e collegato dalla rete locale (\texttt{10.0.1.0/24}). Rappresenta l'utente a rischio minimo.
    \item \textbf{Bob}: operatore interno che utilizza un dispositivo personale privo di attestazione hardware TPM. Pur collegandosi dalla rete aziendale, la mancanza di garanzie sull'integrità del dispositivo lo sottopone a soglie di rischio più restrittive. Visto da un altro punto di vista, Bob rappresenta un esempio di \textbf{insider threat} potenziale, per esempio un utente che ha rubato le credenziali ma non possiede TPM.
    \item \textbf{Charlie}: professionista in regime di smart working, dotato di certificato e chip TPM ma collegato da una rete esterna non controllata (\texttt{192.168.100.0/24}). La provenienza da una rete non fidata impone politiche di tolleranza al rischio più restrittive.
\end{itemize}
Questo scenario consente di dimostrare concretamente come l'architettura adatti dinamicamente le proprie soglie di sicurezza in funzione del contesto multidimensionale di ciascun accesso, combinando l'identità dell'utente, la postura del dispositivo e la localizzazione di rete.

\subsection{Obiettivi del progetto}
Il presente progetto si pone l'obiettivo di realizzare un'architettura Zero Trust completa, funzionante ed integrata per la protezione di un database MongoDB aziendale. Il sistema è progettato per operare una valutazione continua e adattiva del rischio (\textit{Continuous Authorization}) basata sul contesto di rete, sul comportamento dell'utente e sulla postura di sicurezza del dispositivo, inclusa rigorosamente l'attestazione hardware tramite modulo TPM.

Per raggiungere questo scopo, vengono impiegate diverse tecnologie integrate. In primo luogo, \texttt{envoy} opera come \textbf{Policy Enforcement Point} (PEP), intercettando ogni singola richiesta HTTP o MongoDB in modo sicuro mediante crittografia mutua TLS (\textit{mTLS}), eseguendo il fingerprinting software (\textit{JA3}) e l'ispezione profonda dei pacchetti a livello applicativo (\textit{L7 DPI}). In secondo luogo, \texttt{open policy agent} (OPA) funge da \textbf{Policy Decision Point} (PDP) decentralizzato, valutando le richieste in tempo reale sulla base delle sei dimensioni Zero Trust (ZTA 6D), rappresentate dai parametri \textit{u} (user), \textit{d} (device), \textit{s} (software), \textit{n} (network), \textit{a} (action), \textit{r} (resource).

L'analisi comportamentale è affidata a \texttt{splunk} Machine Learning Toolkit (MLTK), interpellato sincronicamente da OPA per calcolare un punteggio di rischio dinamico tramite modelli predittivi basati sull'algoritmo \textbf{Gradient Boosting}. Parallelamente, \texttt{splunk} SIEM si occupa dell'ingestion asincrona dei log per il monitoraggio e la correlazione storica degli eventi. Infine, a livello di rete, \texttt{nftables} e \texttt{snort} (NIDS) garantiscono rispettivamente la microsegmentazione del traffico a livello L3/L4 e il rilevamento tempestivo di tentativi di movimento laterale o bypass del proxy.

\subsection{Struttura della relazione}
Il documento è organizzato come segue: il Capitolo \ref{sec:architettura} illustra l'architettura complessiva della topologia di rete a quattro zone. Il Capitolo \ref{sec:implementazione} scende nel dettaglio dell'implementazione tecnica di ciascun componente di sicurezza, descrivendo il funzionamento di \texttt{envoy}, OPA, \texttt{splunk}, \texttt{snort} e \texttt{nftables}. Il Capitolo \ref{sec:scenari} descrive la validazione del sistema attraverso la simulazione di scenari d'attacco ed accessi legittimi. Infine, il Capitolo \ref{sec:conclusioni} riassume i risultati e propone futuri sviluppi progettuali.
